\documentclass[11pt,a4paper]{report}

\usepackage{xcolor}
\def\farbe{cyan}

\usepackage{dclecture}

\usepackage{pigpen}
\pdfgentounicode=0

\usepackage{morse}
\def\lm{\Large\morse}

\setlength{\headheight}{24pt}



%%% Fancy Header and Footer
\renewcommand{\headrule}{\vbox to 0pt{\hbox to\headwidth{\color{\farbe}\rule{\headwidth}{1pt}}\vss}}
\pagestyle{fancy} %eigener Seitenstil
\fancyhf{} %alle Kopf- und Fusszeilenfelder bereinigen
\fancyhead[C]{Thema} %Kopfzeile mitte
%\fancyhead[R]{\includegraphics[width=0.2cm]{x.png}}
\fancyfoot[C]{\thepage}

\begin{document}

\section*{Mord im Mathematik Department}
Ein Mord wurde im Mathematik-Department an einer
nahegelegenen Schule begangen.

Ihre Aufgabe ist es, die Hinweise zu entschlüsseln, um
\begin{enumerate}
    \item die Identität des Mörders
    \item die Mordwaffe
   \item den Raum, in dem der Mord stattgefunden hat

\end{enumerate}

Die sieben Verdächtigen sind:
\begin{multicols}{4}
    \begin{itemize}
        \item Herr Broome 
        \item Dr. France 
        \item Herr Halai 
        \item Fräulein Kaur
        \item Herr Scotland 
        \item Herr Sheppard 
        \item Herr Smith
    \end{itemize}
\end{multicols}
Die möglichen Mordwaffen sind:
\begin{multicols}{4}
    \begin{itemize}
        \item Meterstab 
        \item Stuhl 
        \item Schere
        \item Zirkel 
        \item Lehrbuch 
        \item Bleistift
        \item Hefter
    \end{itemize}
\end{multicols}
Der Raum, in dem der Mord begangen wurde, könnte sein:
\begin{verbatim}
    16 23 24 25 28 29 Computerraum
\end{verbatim}

\newpage

\section*{Code 1: Schweinepen}
Der erste Code verwendet den folgenden Schlüssel. Sie müssen herausfinden, wie Sie den Schlüssel
verwenden, um die folgende Nachricht zu entschlüsseln.

\centfig{0.8}{pigpen.png}

{\pigpenfont DER RAUM, IN DEM DER MORD BEGANGEN WURDE, HAT EINE RAUMNUMMER}

\section*{Code 2: Polybius-Quadrat}
\centfig{0.5}{polybius}

\begin{verbatim}
    (5,2)(3,4)(5,5) (3,3)(1,1)(3,2)(4,5)(5,5)(3,2)(5,5)(3,2) 
    (4,5)(5,3)(5,5)(4,2) (4,3)(5,3)(5,2) (3,4)(1,5)(2,1)(5,5) 
    (1,5)(4,3) (5,5) (4,4)(4,3) (5,2)(3,4)(5,5)(4,4)(3,2) 
    (4,3)(1,5)(3,3)(5,5)
\end{verbatim}

\section*{Code 3}
\begin{verbatim}
    20,8,5  13,21,18,4,5,18  23,1,19  14,15,20  9,14 1 
    16,18,9,13,5  14,21,13,2,5,18,5,4  18,15,15,13
\end{verbatim}

\section*{Code 4}
Ihr einziger Hinweis für diesen ist, dass er per Mobiltelefon empfangen wurde.

\centfig{0.4}{mobile}

\begin{verbatim}
    843068733709327660363706680428302074277063825072780
\end{verbatim}

\section*{Code 5: Atbash}
Hinweis: Jeder Buchstabe stellt einen anderen Buchstaben in einem einfachen System dar.

\begin{verbatim}
    GSV ILLN MFNYVI RH Z NFOGRKOV LU ULFI
\end{verbatim}

\section*{Code 6: Caesar-Verschlüsselung}

\begin{verbatim}
    KYV DLIUVIVIJ ERDV NZCC KVCC PFL NYRK TFLEKIP YV ZJ WIFD
\end{verbatim}

\section*{Code 7: Morsecode}
\centfig{0.8}{morse}

{\lm die Raum Nummer}

{\lm hat acht Faktoren}

\section*{Code 8: Die letzte Herausforderung}

{\lm kdc qxf mrm qn mx } 

{\lm rc ynaqjyb frcq }

{\lm bxvncqrwp cqjc}

{\lm  bcdmnwcb ljw brc xw}

\end{document}
